\subsection{Validation across image-classification datasets}
\label{sec:dataset_validation}

We next evaluated FROST on MNIST, Fashion-MNIST and CIFAR-10, covering
handwritten digits, grayscale clothing images and RGB natural images
\cite{lecun1998gradient,xiao2017fashion,krizhevsky2009learning}.
Each dataset used a four-block patch Transformer with 16 tokens, two
attention heads per block and a feature width of 32. Images retained their
native resolution and were divided into a $4\times4$ patch grid. Three
reference models were trained independently for each dataset. Their
weights were then frozen, and a shared, rank-8 HI-SMURF operator was
calibrated for each model using separate images from the training
partition. Calibration used reference softmax outputs without task labels.

On the complete official test sets of 10,000 images per dataset,
floating-point accuracies were 97.96\%, 85.50\% and 64.22\%,
respectively. Finite-stream FROST at $L=256$ achieved
96.85\%, 83.95\% and 63.23\%, corresponding to paired accuracy
changes of -1.11, -1.55 and -0.99 percentage points
(Fig.~\ref{fig:dataset_accuracy}). Agreement with the corresponding
floating-point class predictions was 97.84\%, 93.56\% and 91.21\%,
and relative centred-logit errors were 0.201, 0.162 and 0.131.
Analytic FROST achieved 97.70\%, 85.09\% and 64.08\%
accuracy, compared with the finite-stream values above.
The simulated W8A8 control achieved 97.98\%, 85.48\% and 64.16\%
accuracy. This control quantized the linear and attention products while
retaining floating-point nonlinear operations. It provides a numerical
precision comparison; no inference about hardware cost follows from
these accuracy measurements.

We assessed stream-length dependence on a prespecified, class-balanced
subset of 1,000 test images per dataset
(Fig.~\ref{fig:dataset_streams}). Increasing $L$ from 64 to 1,024
changed relative centred-logit errors from 0.383, 0.312 and 0.263
to 0.133, 0.114 and 0.081, respectively. Corresponding accuracies changed
from 91.67\%, 76.70\% and 60.70\% to
97.50\%, 84.00\% and 64.23\%. Analytic evaluation of the same
calibrated finite-state operators yielded errors of
0.105, 0.091 and 0.061 on these subsets. Reporting both classification
accuracy and continuous-output error distinguishes changes in predicted
classes from deviations that leave the classification decision unchanged.

Activation calibration was evaluated against direct transfer of the
generic HI-SMURF operator on the same diagnostic images at $L=256$
(Fig.~\ref{fig:dataset_calibration}). Relative centred-logit errors
changed from 0.287, 0.224 and 0.160 under direct transfer to
0.194, 0.162 and 0.128 after calibration. Prediction agreement changed
from 95.23\%, 89.40\% and 89.53\% to
98.03\%, 92.80\% and 91.77\%. These comparisons preserve the
Transformer weights, FSM state count and TT rank. They assess adaptation
of a fixed-capacity operator to each model's activation distribution;
they do not establish zero-shot transfer of one calibrated operator
between datasets.

Reported central values are means across three independently trained
models; figure error bars denote their sample standard deviations.
Each complete-test evaluation uses one prespecified stream realization.
Additional realizations at $L=256$ were evaluated on the diagnostic
subsets and are reported separately. Together, these experiments quantify
task-level performance and numerical fidelity of the complete stochastic
encoder across three image datasets, while retaining the limitations of
a compact backbone trained from scratch. The simulation includes
numerical centering, pooling, residual and bias addition, range handling
and stage conversions. Hardware area, energy and latency were not
evaluated.
